%% kapa: exam=1 round=35 year=2024 period=1 form=A subject=economics subject_ko=경제학원론 seq=28 number=68
%% source: Q-Net 감정평가사 기출문제 (cst00309) artlSeq=5213229 file=raw/1008/5213229/1교시 문제지 원본.pdf
%% conversion: kordoc 4.24.51 -> tools/gen1.py (auto); review=proofread-sonnet; flags=none
\begin{question}{68}
소비이론에 관한 설명 중 옳은 것은?
\begin{choices}
  \item 케인즈(Keynes)의 소비이론에 따르면 이자율이 소비의 주요 결정요인이다.
  \item 생애주기가설에 따르면 은퇴연령의 변화 없이 기대수명이 증가하면 소비가 감소한다.
  \item 리카도 등가(Ricardian equivalence)정리는 케인즈의 소비함수에 기초한 이론이다.
  \item 케인즈의 소비이론은 소비자들의 소비평탄화(consumption smoothing)를 강조한다.
  \item 소비에 대한 임의보행(random walk)가설은 유동성제약에 직면한 소비자의 소비 선택을 설명한다.
\end{choices}
\end{question}
\answer{2}
